\documentclass{article}
\usepackage{graphicx} % Required for inserting images
\usepackage[heading,UTF8]{ctex}             % 中文支持
\usepackage{geometry}                                       % 页面设置
\usepackage{fancyhdr,lastpage}                              % 页眉页脚
\usepackage{fontspec}                                       % 字体
\usepackage{tabularx,multirow,booktabs}                     % 表格
\usepackage{multicol}                                       % 分栏
\usepackage{listings}                                       % 代码
\usepackage{algorithm,algorithmicx,algpseudocode}           % 伪代码
\usepackage{amsmath,amsthm,amsfonts,amssymb,bm}             % 数学公式
\usepackage{epigraph}                                       % 引用
\usepackage{hyperref}                                       % 超链接
\usepackage{tikz}                                           % TikZ 画图
\usepackage{tcolorbox}                                      % Color Box
\usepackage{enumitem}                                       % 列表样式
\usepackage{xeCJKfntef,xcolor}                              % 文字加点
\usepackage{diagbox}



\hypersetup{                                                % 超链接样式
    colorlinks=true,
    linkcolor=blue,
    filecolor=blue,
    urlcolor=blue,
    citecolor=cyan,
}

\usetikzlibrary{arrows.meta}                                % 附加箭头样式

\tikzset{                                                   % TikZ 设置
    global scale/.style={                                     % 整张图放缩
        scale=#1,
        every node/.append style={scale=#1}
    }
}

\tcbset{                                                    % Color Box 样式
    colback=white,
    colframe=black,
    boxrule=0.5pt
}

% \lstset{                                                    % 代码框格式设置
%     breaklines = true,
%     keepspaces = true,
%     basicstyle = \ttfamily,
%     numberstyle = \footnotesize\ttfamily\color{gray},
%     numbers = left,
%     columns = flexible,
%     framexleftmargin = 0pt,
%     xleftmargin = 12pt,
%     numbersep = 10pt,
%     keywordstyle = \color{black}, %\color{blue!70},
%     commentstyle = \color{black}, %\color{green!70},
%     frame = single,
%     frameround = tttt,
%     rulecolor = \color{blue},
%     language = c++,
%     escapeinside = {(*@}{@*)}
% }

% 样例的代码框格式设置 
\lstdefinestyle{sample} {
    breaklines = true,
    keepspaces = true,
    basicstyle = \ttfamily,
    numberstyle = \footnotesize\ttfamily\color{gray},
    numbers = left,
    columns = flexible,
    framexleftmargin = 0pt,
    xleftmargin = 12pt,
    numbersep = 10pt,
    keywordstyle = \color{black}, %\color{blue!70},
    commentstyle = \color{black}, %\color{green!70},
    frame = single,
    frameround = tttt,
    rulecolor = \color{blue},
    language = c++,
    escapeinside = {(*@}{@*)}
}

\definecolor{mygreen}{RGB}{0,128,0}
% 题解代码代码框格式设置
\lstdefinestyle{code}{
    breaklines = true,
    keepspaces = true,
    basicstyle = \ttfamily,
    numberstyle = \footnotesize\ttfamily\color{gray},
    keywordstyle = \color{blue!70}, 
    commentstyle = \color{gray}, 
    stringstyle = \color{mygreen},
    numbers = left,
    columns = flexible,
    framexleftmargin = 0pt,
    xleftmargin = 12pt,
    numbersep = 10pt,
    frame = single,
    frameround = tttt,
    rulecolor = \color{blue},
    language = c++,
    escapeinside = {(*@}{@*)},
    showstringspaces = false
}
\xeCJKsetup{                                                % 文字下方加点样式
    underdot/symbol = {\color{black}.}
}

\newcommand{\filename}[1]{{\textbf{\textit{#1}}}}           % 自定义文件名样式
\newcommand{\myemph}[1]{{\textbf{\CJKunderdot{#1}}}}        % 自定义强调样式
\ctexset{                                                   % 标题设置
    section = {
        format = \centering\Large\heiti,
        numbering = true,
        number = {},
        aftername = {}
    },
    subsection = {
        format = \fontsize{13pt}{\baselineskip}\heiti,
        numbering = true,
        number = {},
        aftername = {【},
        aftertitle = {】}
    },
    subsubsection = {
        format = \myemph,
        numbering = true,
        number = {},
        aftername = {}
    }
}

\geometry{                                                  % 页面设置
    a4paper,
    left = 70pt,
    right = 70pt,
    top = 80pt,
    bottom = 65pt,
    footskip = 30pt
}

\floatname{algorithm}{算法}
\renewcommand{\algorithmicrequire}{\textbf{输入：}}
\renewcommand{\algorithmicensure}{\textbf{输出：}}

\setmonofont{Consolas} [
    Path = ./, 
    Extension = .ttf,
]              % 等宽字体


\setlist[itemize]{itemsep=0pt,partopsep=0pt,parsep=0pt,topsep=10pt}
\setlist[enumerate]{itemsep=0pt,partopsep=0pt,parsep=0pt,topsep=10pt}
\setlist[itemize,1]{leftmargin=35pt}
\setlist[enumerate,1]{leftmargin=35pt}

\newcommand{\docTitle}{CSP-S算法模版文档}
\newcommand{\docSecTitle}{}
\newcommand{\docThiTitle}{}
\newcommand{\docAuthor}{hezlik}
\newcommand{\docDate}{年月日 $\sim$ }

\newcommand{\fancyheadLeftText}{\docTitle}
\newcommand{\fancyheadRightText}{\docThiTitle}

\title{\docTitle}
\author{\docAuthor}
\date{\docDate}

\hypersetup{
    pdftitle={\docTitle},
    pdfauthor={\docAuthor}
}

\makeatletter
\renewcommand*\maketitle{                                   % 标题页格式设置
    \begin{center}
        {\LARGE \bfseries \@title \par}
        {\Huge \heiti \docSecTitle \par}
        {\LARGE \kaishu \docThiTitle \par}
        {\large \@author \par}
        {\large \heiti 时间：\@date \par}
    \end{center}
    \setcounter{footnote}{0}
}
\makeatother


\begin{document}
    \renewcommand\headrulewidth{0.5pt}
    \pagestyle{fancy}
    \fancyhf{}
    \fancyhead[L]{\footnotesize \fancyheadLeftText}
    \fancyfoot[C]{\footnotesize 第 \thepage 页\quad 共 \pageref{LastPage} 页}
    
    \section{单源最短路$\&$判断负环(Bellman-Ford+SPFA)} % 大标题
    \fancyhead[R]{\footnotesize \fancyheadRightText~单源最短路$\&$判断负环} % 页眉
    
    \subsection{应用范围}
    
    单源最短路指的是求解两点之间的最短距离，通常出现在图论题目当中。

    常见的单源最短路算法有$Floyd$,$Bellman-Ford$,$Dijkstra$,本文主要介绍$Bellman-Ford$算法以及它的队列优化版本$SPFA$。

    
    \subsection{算法思路}
    1、定义数组$Dist[i]$。$Dist[i]$代表了从起点到达点$i$的距离，并将它们初始化成$Inf$，表示不可达。
    
    2、定义一个操作：松弛操作。选取一个边$(u , v)$，利用到达$u$的距离来更新到达$v$的距离，即$min(Dist[v] , Dist[u] + len(u,v))$。

    3、$Bellman-Ford$要做的就是每一轮都把每一条边松弛一遍，直到一轮中没有任何一个距离被更新为止。
 \begin{lstlisting}[style = code]
 while (flag) { 
    // 枚举每个点
    flag = false;
    for (int i = 1; i <= n; i++) {
        for (int j = 0; j < e[i].size(); j++) {
            int nxt = e[i][j].to;
            int val = e[i][j].val;
            // i -> nxt: val;
            if (dist[i] + val < dist[nxt]) {
                flag = true;
                dist[nxt] = min(dist[nxt], dist[i] + val);
            }
        }
    }
}\end{lstlisting}

    4、由于一条简单路径最多只会经过$n-1$条边，故最多进行$n-1$轮，整体复杂度为$O(nm)$。

    5、特别的，但如果从 $S$ 点出发，抵达一个负环时，松驰操作会无休止地进行下去。注意到前面的论证中已经说明了，对于最短路径存在的图，松驰操作最多只会执行 $n-1$ 轮，因此如果第 $n$ 轮循环时仍然存在能松驰的边，说明从 $S$ 点出发，能够抵达一个负环。

    6、在$Bellman-Ford$算法中，我们每一轮都是把每条边都尝试松弛，接下来考虑从减少每一轮要松弛的边入手，优化复杂度。
\newpage
    7、我们会发现：只有在一轮中被松弛的点所延申出的边，才有可能在下一轮引发松弛操作。因此我们考虑用队列来维护\textbf{可能引发松弛操作的点}，就可以减少每一轮所需要松弛的边了。这就是$SPFA$算法的核心。
     \begin{lstlisting}[style = code]
    queue<int> q;
    q.push(s);
    dis[s] = 0;
    while (!q.empty()) {
        int now = q.front();
        q.pop();
        for (int i = 0; i < e[now].size(); i++) {
            int nxt = e[now][i].to;
            int val = e[now][i].val;
            if (dis[now] + val < dis[nxt] ) {
                q.push(nxt);
                dis[nxt] = dis[now] + val;
            }
        }
    }\end{lstlisting}
    
    8、$SPFA$同样也可以判断负环，只需记录最短路经过了多少条边，当经过了至少$n$条边时，说明$s$点可以抵达一个负环。
    \newpage
    \subsection{样例题意}
输入的第一行是一个整数 $T$，表示测试数据的组数。对于每组数据的格式如下：

第一行有两个整数，分别表示图的点数 $n$ 和接下来给出边信息的条数 $m$。

接下来 $m$ 行，每行三个整数 $u$, $v$, $w$。
 \subsection{样例 1 输入} % 样例使用style = sample

\begin{lstlisting}[style = sample]
2
3 4
1 2 2
1 3 4
2 3 1
3 1 -3
3 3
1 2 3
2 3 4
3 1 -8
\end{lstlisting}
    
    
    \subsection{样例 1 输出}
    
\begin{lstlisting}[style = sample]
NO
YES
\end{lstlisting}
    

    \newpage
    \subsection{代码模版} 

    \begin{lstlisting}[style = code]
#include<bits/stdc++.h>
using namespace std;
int T, n, m, u, v, w;
const int maxn = 1e5 + 10;
struct node {
    int to, val;
};
vector <node>  e[maxn];
int dis[maxn];
int cnt[maxn];//最短方案中，到达点i的边的个数
bool in[maxn];
void spfa(){
    memset(dis, 0x3f3f, sizeof(dis));
    memset(in, 0, sizeof(in));
    memset(cnt, 0, sizeof(cnt));
    queue <int> q;
    q.push(1);
    dis[1] = 0;
    in[1] = true;
    cnt[1] = 1;
    while (!q.empty()){
        u = q.front();
        q.pop();
        in[u] = false;
        for (int i = 0; i < e[u].size(); i++){
            v = e[u][i].to;
            w = e[u][i].val;
            if (dis[v] > dis[u] + w){
                dis[v] = dis[u] + w;
                if (!in[v]){
                    cnt[v]++;
                    q.push(v);
                    in[v] = true;
                    if (cnt[v] >= n){
                        printf("YES\n");
                        return;
                    }
                }
            }
        }
    }
    cout << "NO\n";
}
int main(){
    cin >> T;
    for (int  t = 0; t < T; t++){
        cin >> n >> m;
        for (int i = 1; i <= n; i++){
            e[i].clear();
        }
        for (int i = 0; i < m; i++){
            cin >> u >> v >> w;
            e[u].push_back({v , w});
            if (w >= 0){
                e[v].push_back({u , w});
            }
        }
        spfa();
    }
    return 0;
}\end{lstlisting}
    
    
    
    

    
    
\end{document}